\section{Methodology}

\subsection{Neuromorphic Reservoir Core and Fixed-Point Engine}
The Eón core processes continuous sequential inputs through a discrete-time, leaky-integrator Echo State Network (ESN) \cite{jaeger2001echo}. Let $u(t) \in \mathbb{R}^K$ denote the external input vector at time step $t$, and let $x(t) \in \mathbb{R}^N$ denote the internal state of the recurrent reservoir consisting of $N$ neurons. The dynamic state update equation is formulated as:
\begin{equation}
    \tilde{x}(t) = \tanh\left( W_{in} u(t) + W x(t-1) + \nu(t) \right)
    \label{eq:state_pre}
\end{equation}
\begin{equation}
    x(t) = (1 - \alpha) x(t-1) + \alpha \tilde{x}(t)
    \label{eq:state_leak}
\end{equation}
where $\alpha \in (0, 1]$ is the leak rate governing the effective temporal integration timescale, $W_{in} \in \mathbb{R}^{N \times K}$ is the input projection matrix, $W \in \mathbb{R}^{N \times N}$ is the recurrent reservoir weight matrix, and $\nu(t) \sim \mathcal{N}(0, \sigma^2)$ is an optional low-amplitude regularization noise term.

The predicted output $\hat{y}(t) \in \mathbb{R}^L$ is computed via a linear projection of the extended state vector $z(t) = [1; u(t); x(t)]^T \in \mathbb{R}^{M}$ (where $M = 1 + K + N$, with the leading $1$ representing the constant bias term):
\begin{equation}
    \hat{y}(t) = W_{out} z(t)
    \label{eq:readout}
\end{equation}
where $W_{out} \in \mathbb{R}^{L \times M}$ is the sole trainable matrix in the architecture.

\paragraph{Fixed-Point Q8.8 Arithmetic:}
To execute on deeply embedded microcontrollers devoid of hardware Floating-Point Units (FPUs), Eón implements a native Q8.8 signed fixed-point representation. A real number $v \in [-128.0, 127.996]$ is represented as a 16-bit signed integer $q = \lfloor v \cdot 2^8 \rceil \in [-32768, 32767]$. 
Addition and subtraction operate natively in standard integer ALU registers. Multiplication of two Q8.8 operands requires an intermediate 32-bit product followed by an arithmetic right-shift of $8$ bits:
\begin{equation}
    q_{prod} = \text{int16}\left( (\text{int32}(q_a) \times \text{int32}(q_b)) \gg 8 \right)
\end{equation}
The hyperbolic tangent non-linearity is evaluated either via a compact 256-entry lookup table (LUT) or an efficient piecewise linear/polynomial approximation:
\begin{equation}
    \tanh(v) \approx \begin{cases}
        v - \frac{1}{3} v^3 & \text{if } |v| < 1.0 \\
        \text{sign}(v) & \text{if } |v| \ge 2.0 \\
        \text{sign}(v) \cdot (0.76 + 0.12 |v|) & \text{otherwise}
    \end{cases}
\end{equation}
limiting the maximum approximation error to under $0.5\%$ while consuming zero dynamic heap allocations.

\subsection{Orthogonal Reservoirs via Cayley Transform}
The classical Echo State Property (ESP) requires that the effect of initial reservoir conditions $x(0)$ asymptotically vanishes over time, which is strictly satisfied when the spectral radius $\rho(W) = \max_i |\lambda_i(W)| < 1$. 

In standard sparse Gaussian reservoirs, eigenvalues are non-uniformly distributed inside the spectral disk according to Girko's circular law. Modes corresponding to eigenvalues with $|\lambda_i| \ll 1$ decay exponentially fast, resulting in a suboptimal memory capacity. To achieve theoretically maximal linear memory capacity ($MC \approx N$), Eón implements an \textbf{Orthogonal Reservoir} generator using the algebraic \textbf{Cayley Transform}. 

Let $A \in \mathbb{R}^{N \times N}$ be a random skew-symmetric matrix such that $A^T = -A$. The Cayley transform maps the real Lie algebra of skew-symmetric matrices into the Lie group of special orthogonal matrices:
\begin{equation}
    Q = (I - A)(I + A)^{-1}
    \label{eq:cayley}
\end{equation}
The scaled reservoir matrix is defined as $W = \rho_{target} \cdot Q$. 
\paragraph{Proposition 1:} \textit{For every eigenvalue $\lambda$ of $W$, $|\lambda| = \rho_{target}$.}
\begin{proof}
Let $v$ be an eigenvector of $A$ with eigenvalue $i \mu$ ($\mu \in \mathbb{R}$, since skew-symmetric matrices possess purely imaginary spectra). Then $(I - A) v = (1 - i \mu) v$ and $(I + A) v = (1 + i \mu) v$. Hence, $Q v = \frac{1 - i \mu}{1 + i \mu} v$. The eigenvalue of $Q$ has magnitude $\left| \frac{1 - i \mu}{1 + i \mu} \right| = \sqrt{\frac{1 + \mu^2}{1 + \mu^2}} = 1$. Scaling by $\rho_{target}$ guarantees that all eigenvalues lie strictly on the circle of radius $\rho_{target}$.
\end{proof}
This formulation preserves energy across all dynamic state dimensions, eliminating dead modes and maximizing echo retention in non-stationary signals.

\subsection{Two-Stage Criticality Screening Filter}
For platforms where on-device seed search or procedural reservoir instantiation is required, calculating full LAPACK eigenspectra ($O(N^3)$ operations) to find seeds with natural edge-of-chaos criticality ($\rho \approx 1.0$) is computationally prohibitive. 

By Girko's circular law, an unscaled random sparse matrix with Gaussian entries $\mathcal{N}(0, 1)$ and density $p$ has an expected spectral radius of:
\begin{equation}
    \mathbb{E}[\rho(W)] \approx \sqrt{N \cdot p}
\end{equation}
For $N=50, p=0.1$, $\mathbb{E}[\rho] \approx 2.24$. Random seeds that naturally fall within the critical resonance window $\rho \in [0.90, 1.10]$ represent less than $2\%$ of candidates. Eón utilizes a \textbf{Two-Stage Power Iteration Filter}:
\begin{enumerate}
    \item \textbf{Stage 1 (Rayleigh Quotient Screening):} Performs $k=8$ power iterations starting from a unit vector $v_0$:
    \begin{equation}
        v_{m+1} = \frac{W v_m}{\|W v_m\|_2}, \qquad \lambda_{approx} = \|W v_k\|_2
    \end{equation}
    If $|\lambda_{approx} - 1.0| > \delta_{threshold}$, the candidate seed is rejected immediately in $O(k \cdot p N^2)$ operations ($<15\,\mu\text{s}$ per trial).
    \item \textbf{Stage 2 (Exact Verification):} Only the surviving candidate seeds ($<2\%$) are subjected to full QR eigenvalue verification.
\end{enumerate}
This two-stage pipeline reduces seed discovery latency from 15 minutes to under 3 seconds on standard edge microprocessors, representing a $>98\%$ compute reduction.

\subsection{Convex Readout Training via Exact Cholesky Factorization}
Given a collection of $T$ training state vectors collected after a washout period $T_w$, let $Z \in \mathbb{R}^{M \times T_{train}}$ denote the state matrix and $Y \in \mathbb{R}^{L \times T_{train}}$ denote the target sequence. The regularized ridge regression objective minimizes:
\begin{equation}
    \mathcal{J}(W_{out}) = \frac{1}{2} \|Y - W_{out} Z\|_F^2 + \frac{\lambda}{2} \|W_{out}\|_F^2
    \label{eq:ridge_loss}
\end{equation}
Setting $\nabla_{W_{out}} \mathcal{J} = 0$ yields the classical normal equations:
\begin{equation}
    (Z Z^T + \lambda I_M) W_{out}^T = Z Y^T
    \label{eq:normal_eq}
\end{equation}
Let $A = Z Z^T + \lambda I_M \in \mathbb{R}^{M \times M}$ and $B = Z Y^T \in \mathbb{R}^{M \times L}$. 

\paragraph{Proposition 2:} \textit{The coefficient matrix $A$ is strictly symmetric positive-definite (SPD) for all $\lambda > 0$.}
\begin{proof}
Symmetry is evident as $(Z Z^T + \lambda I)^T = (Z^T)^T Z^T + \lambda I = Z Z^T + \lambda I$. For any non-zero vector $v \in \mathbb{R}^M \setminus \{0\}$:
\begin{equation}
    v^T A v = v^T (Z Z^T) v + \lambda v^T I v = \|Z^T v\|_2^2 + \lambda \|v\|_2^2
\end{equation}
Since $\|Z^T v\|_2^2 \ge 0$ and $\lambda > 0$, $v^T A v \ge \lambda \|v\|_2^2 > 0$. Thus, all eigenvalues of $A$ satisfy $\sigma_i(A) \ge \lambda > 0$, guaranteeing strict positive-definiteness and a condition number bounded by $\kappa(A) \le 1 + \sigma_{max}(Z Z^T) / \lambda$.
\end{proof}

Exploiting the SPD property, Eón computes the unique \textbf{Cholesky Factorization}:
\begin{equation}
    A = L \cdot L^T
    \label{eq:cholesky}
\end{equation}
where $L$ is a real lower triangular matrix with strictly positive diagonal entries:
\begin{equation}
    L_{j,j} = \sqrt{A_{j,j} - \sum_{k=0}^{j-1} L_{j,k}^2}
\end{equation}
\begin{equation}
    L_{i,j} = \frac{1}{L_{j,j}} \left( A_{i,j} - \sum_{k=0}^{j-1} L_{i,k} L_{j,k} \right), \quad \text{for } i > j
\end{equation}
The optimal weight matrix is then resolved via two triangular substitution passes:
\begin{enumerate}
    \item \textbf{Forward Substitution:} Solve $L U = B$ for intermediate matrix $U \in \mathbb{R}^{M \times L}$:
    \begin{equation}
        U_{i, \ell} = \frac{1}{L_{i,i}} \left( B_{i, \ell} - \sum_{k=0}^{i-1} L_{i,k} U_{k, \ell} \right)
    \end{equation}
    \item \textbf{Backward Substitution:} Solve $L^T W_{out}^T = U$ for $W_{out}^T$:
    \begin{equation}
        (W_{out}^T)_{i, \ell} = \frac{1}{L_{i,i}} \left( U_{i, \ell} - \sum_{k=i+1}^{M-1} L_{k,i} (W_{out}^T)_{k, \ell} \right)
    \end{equation}
\end{enumerate}

\begin{table}[t]
\centering
\caption{Computational Complexity and Stability of Ridge Solvers}
\label{tab:solver_complexity}
\begin{tabular}{lccc}
\toprule
\textbf{Algorithm} & \textbf{Flops ($M^3$)} & \textbf{Stability} & \textbf{Pivoting} \\
\midrule
Gauss-Jordan Inversion & $1.00\, M^3$ & Poor & Required \\
LU Decomposition ($PA = LU$) & $0.67\, M^3$ & Moderate & Required \\
\textbf{Cholesky Factorization ($LL^T$)} & $\mathbf{0.33\, M^3}$ & \textbf{Inconditionally Stable} & \textbf{None} \\
\bottomrule
\end{tabular}
\end{table}

As shown in Table~\ref{tab:solver_complexity}, Cholesky factorization cuts the arithmetic complexity by $66\%$ compared to Gauss-Jordan elimination and $50\%$ compared to general LU factorization. Crucially, because $L$ requires no row pivoting, the memory access pattern is completely deterministic, avoiding branching overhead and cache misses on microcontrollers.

\subsection{Online Normalized LMS (NLMS) Adaptation}
While batch Cholesky factorization provides optimal offline training, resource-constrained sensors deployed in the wild cannot afford to store historical activation matrices $Z \in \mathbb{R}^{M \times T_{train}}$ in RAM. Furthermore, Recursive Least Squares (RLS) requires maintaining an $M \times M$ inverse covariance matrix $P(t)$, requiring $O(M^2)$ memory and operations.

To enable true continuous on-device learning in under $2\,\text{KB}$ RAM, Eón introduces an embedded \textbf{Normalized Least-Mean-Squares (NLMS)} adaptation rule in \texttt{aeon\_online\_train\_step}:
\begin{equation}
    e(t) = y_{target}(t) - \hat{y}(t) = y_{target}(t) - W_{out}(t) z(t)
\end{equation}
\begin{equation}
    W_{out}(t+1) = W_{out}(t) + \eta \cdot \frac{e(t) \cdot z(t)^T}{\|z(t)\|_2^2 + \epsilon}
    \label{eq:nlms}
\end{equation}
where $\eta \in (0, 1]$ is the adaptive learning rate, and $\epsilon = 10^{-6}$ is a small regularization constant preventing division by zero.
\paragraph{Properties of the Embedded NLMS Rule:}
\begin{itemize}
    \item \textbf{Memory Overhead:} Exactly $O(1)$ auxiliary RAM; no covariance matrix or historical activation buffer is retained.
    \item \textbf{Arithmetic Complexity:} $2M$ MAC operations per sample, executing in $<2\,\mu\text{s}$ on an $80\,\text{MHz}$ Cortex-M4F.
    \item \textbf{Energy Scale Invariance:} Normalizing by the instantaneous power $\|z(t)\|_2^2$ ensures numerical stability regardless of input transducer amplitude scaling.
\end{itemize}

\subsection{Asymptotically Bounded Synaptic Plasticity (Oja's Rule)}
When unsupervised continuous adaptation of the internal reservoir $W$ is required to track non-stationary environmental dynamics, classical Hebbian learning ($\Delta W = \eta x(t) x(t-1)^T$) leads to exponential weight explosion. To prevent destabilization of the Echo State Property without running costly per-step spectral renormalizations, Eón implements \textbf{Oja's Normalized Plasticity Rule} \cite{oja1982simplified}:
\begin{equation}
    \Delta w_{i,j} = \eta \cdot \left( x_i(t) \cdot x_j(t-1) - x_i(t)^2 \cdot w_{i,j} \right)
    \label{eq:oja}
\end{equation}
In compact matrix notation:
\begin{equation}
    \Delta W = \eta \cdot \left( x(t) x(t-1)^T - \text{diag}(x(t)^2) W \right)
\end{equation}
The negative damping term $-x_i(t)^2 w_{i,j}$ provides intrinsic homeostatic weight regularization. As proven by Oja \cite{oja1982simplified}, this rule asymptotically constrains synaptic vectors to the unit sphere ($\|w_i\|_2 \to 1$) and converges to the principal eigenvector subspace of the input autocorrelation matrix. Eón restricts spectral radius re-verification to once every 10 steps, achieving a $10\times$ speedup during continuous online tracking.

\subsection{Deterministic Sync \& 1-Bit Weight Quantization Protocol}
In distributed multi-agent edge deployments (e.g., wireless sensor swarms), exchanging full floating-point weight matrices is prohibited by radio duty cycles and power limits. Eón resolves this through two complementary mechanisms:
\begin{enumerate}
    \item \textbf{Zero-Bandwidth Procedural Synchronization (Spirit Hash):} Nodes agree on a 32-bit root seed $S_0$ derived from a synchronized network timestamp and device identity. Weights are deterministically generated on-demand using a shared Linear Congruential Generator (LCG):
    \begin{equation}
        S_{k+1} = (1103515245 \cdot S_k + 12345) \bmod 2^{31}
    \end{equation}
    This ensures that two physically isolated microcontrollers instantiate structurally identical reservoir matrices without transmitting a single weight over the air.
    \item \textbf{1-Bit Sign Quantization:} When nodes collaborate by sharing adapted readout vectors $W_{out}$, weights are quantized to their 1-bit sign representation:
    \begin{equation}
        b_i = \begin{cases}
            1 & \text{if } W_{out}[i] \ge 0 \\
            0 & \text{if } W_{out}[i] < 0
        \end{cases}
    \end{equation}
    For an $N=50$ reservoir, the readout payload condenses to exactly $\lceil 52 / 8 \rceil = 7$ bytes. Including a 14-byte cryptographic header, the complete transmission packet occupies only $21$ bytes, compared to $175$ bytes for JSON and $214$ bytes for binary float32.
\end{enumerate}
